% ============================================================================
% COURS 03 — MODÉLISATION DES TRANSMETTEURS — PARTIE A
% Source : 4-ModelisationTransmetteurs2022.docx
% ============================================================================

\section{Modélisation des transmetteurs}
\label{sec:trans-objectifs}

\begin{TLThreeMethod}{Objectifs de la partie transmetteurs}
Déterminer les lois entrée-sortie et les rapports de transmission des
transmetteurs usuels, distinguer réduction et multiplication, prendre en
compte les sens de rotation, la transformation de mouvement, le
rendement et la réversibilité.
\end{TLThreeMethod}

\begin{center}
\small
\begin{tabularx}{\linewidth}{>{\bfseries}p{2cm}X c}
\toprule
Je connais & Connaissances attendues & \\
\midrule
& Les symboles et les familles de transmetteurs usuels. & $\square$\\
& Le rapport de transmission d'un engrenage, d'un train simple, d'un train épicycloïdal et d'une roue-vis sans fin. & $\square$\\
& Les lois d'un pignon-crémaillère, d'une vis-écrou et d'un dispositif poulies-courroie ou pignons-chaîne. & $\square$\\
\midrule
Je sais & Identifier les roues menantes et menées et compter les contacts extérieurs. & $\square$\\
& Déterminer un rapport global puis relier vitesse, couple et rendement. & $\square$\\
& Écrire la formule de Willis et exploiter une configuration de train épicycloïdal. & $\square$\\
\bottomrule
\end{tabularx}
\end{center}

\subsection{Mise en situation : objectif photographique}

L'autofocus déplace une lentille pour faire converger les rayons lumineux
sur le capteur. Une machine à courant continu, plusieurs étages de
transmission et un mécanisme vis-écrou réalisent ce déplacement.

\begin{center}
\TLFourImage[.31\linewidth]{image5.png}\hfill
\TLFourImage[.31\linewidth]{image6.png}\hfill
\TLFourImage[.31\linewidth]{image7.png}
\end{center}

Pour vérifier la résolution annoncée par le constructeur, il faut relier
la rotation du moteur au déplacement de la lentille, puis tenir compte
des codeurs installés dans la chaîne.

\section{Grandeurs d'entrée et de sortie}
\label{sec:trans-grandeurs}

Pour un transmetteur, l'entrée et la sortie sont définies par la chaîne
de puissance. Le mouvement d'entrée est souvent continu ; le mouvement
de sortie peut être continu, alternatif ou intermittent.

\begin{TLThreeDef}{Loi entrée-sortie d'un transmetteur}
Une loi entrée-sortie est une relation imposée par la structure du
transmetteur entre les grandeurs d'entrée et de sortie :
\[
\omega_s=f(\omega_e),
\qquad
v_s=f(\omega_e).
\]
Pour une transmission rotation--rotation linéaire :
\[
\boxed{r=\frac{\omega_s}{\omega_e}
=\frac{\dot\theta_s}{\dot\theta_e}.}
\]
Pour une transformation rotation--translation :
\[
\boxed{k=\frac{v_s}{\omega_e}.}
\]
Le coefficient $r$ est sans dimension ; $k$ s'exprime en mètres.
\end{TLThreeDef}

Lorsque l'entrée et la sortie peuvent être permutées, le transmetteur est
\textbf{réversible}. Dans le cas contraire, il est irréversible.

\subsection{Vitesse, couple, puissance et rendement}

Avec $P=C\omega$ :
\[
\eta=\frac{P_s}{P_e}
=\frac{C_s\omega_s}{C_e\omega_e}
=\frac{C_s}{C_e}\,r.
\]
Ainsi :
\[
\boxed{\frac{C_s}{C_e}=\frac{\eta}{r}.}
\]
Un réducteur de vitesse, $|r|<1$, multiplie généralement le couple. Un
multiplicateur, $|r|>1$, réduit le couple disponible.

Rappel de conversion :
\[
\boxed{\omega=\frac{2\pi N}{60}}
\qquad
(N\text{ en tr/min},\ \omega\text{ en rad/s}).
\]

\section{Classification des transmissions de puissance}
\label{sec:trans-classification}

\begin{center}
\small
\begin{tabularx}{\linewidth}{>{\bfseries}p{3cm}p{3cm}X}
\toprule
Famille & Mouvements & Exemples \\
\midrule
Sans transformation & rotation $\rightarrow$ rotation & roues de friction, engrenages, trains simples, roue-vis sans fin, poulies-courroie, pignons-chaîne \\
Avec transformation & rotation $\leftrightarrow$ translation & pignon-crémaillère, vis-écrou, cylindre roulant, courroie utilisée comme organe linéaire \\
Mouvement alternatif & rotation continue $\rightarrow$ translation ou rotation alternative & bielle-manivelle, excentrique, came \\
Mouvement intermittent & rotation continue $\rightarrow$ rotation discontinue & croix de Malte \\
\bottomrule
\end{tabularx}
\end{center}

Une transformation de mouvement peut modifier la nature, la direction,
le sens ou la vitesse du mouvement.

\section{Transmetteurs sans transformation de mouvement}
\label{sec:trans-sans}

\subsection{Roues de friction}

Deux roues cylindriques sont maintenues en contact par un effort presseur.
L'adhérence doit être suffisante pour empêcher le glissement.

\begin{center}\TLRouesFriction\end{center}

Au point de contact $I$, le roulement sans glissement impose :
\[
\vec V(I,1/2)=\vec0.
\]
Pour deux roues en contact extérieur :
\[
\boxed{\frac{\omega_2}{\omega_1}=-\frac{R_1}{R_2}
=-\frac{D_1}{D_2}.}
\]
Les roues tournent en sens opposés. La puissance transmissible est
limitée par le coefficient d'adhérence et l'effort presseur. Ce principe
est utilisé pour une dynamo de vélo, un home-trainer ou certains anciens
tourne-disques.

\subsection{Engrenage simple}

Un engrenage associe deux roues dentées. La denture impose la même
cinématique que deux roues de friction sur leurs cercles primitifs.

\begin{center}\TLEngrenageSimple\end{center}

\begin{TLThreeDef}{Géométrie de la denture}
Pour une roue de $Z$ dents et de module $m$ :
\[
\boxed{d=mZ,}
\qquad
\boxed{p=\frac{\pi d}{Z}=\pi m.}
\]
Deux roues ne peuvent engrener que si elles possèdent le même module et
le même pas primitif. L'angle de pression normalisé le plus courant vaut
$20^\circ$.
\end{TLThreeDef}

Pour un engrenage extérieur :
\[
\boxed{\frac{\omega_2}{\omega_1}=-\frac{Z_1}{Z_2}.}
\]
Pour un engrenage intérieur, les deux roues tournent dans le même sens :
\[
\boxed{\frac{\omega_2}{\omega_1}=+\frac{Z_1}{Z_2}.}
\]

L'effort de denture se décompose en :
\begin{itemize}
  \item une composante tangentielle $F_t$, seule à transmettre la puissance ;
  \item une composante radiale $F_r$, sollicitant les arbres en flexion ;
  \item une composante axiale $F_a$ pour les dentures hélicoïdales ou coniques.
\end{itemize}

\subsection{Train d'engrenages simple}

Dans un train simple, tous les axes de rotation sont fixes par rapport au
bâti. Le rapport global est le produit des rapports élémentaires :
\[
\boxed{r=\frac{\omega_s}{\omega_e}
=(-1)^n\frac{\prod Z_{\text{menantes}}}
{\prod Z_{\text{menées}}}.}
\]
Le nombre $n$ est le nombre de contacts extérieurs. Le signe n'a de sens
que si les axes d'entrée et de sortie sont parallèles.

Une roue à la fois menée et menante est une \textbf{roue folle} : son
nombre de dents se simplifie dans le rapport, mais elle peut modifier le
sens de rotation.

\begin{TLThreeMethod}{Calculer le rapport d'un train simple}
\begin{enumerate}
  \item Identifier l'entrée, la sortie et les roues solidaires d'un même arbre.
  \item Pour chaque engrènement, repérer la roue menante et la roue menée.
  \item Multiplier les nombres de dents des menantes et des menées.
  \item Compter les contacts extérieurs pour déterminer le signe.
  \item Vérifier que $|r|<1$ pour un réducteur et $|r|>1$ pour un multiplicateur.
\end{enumerate}
\end{TLThreeMethod}

\begin{TLThreeApp}{Exemple de train simple}
Pour un train à trois étages :
\[
r=\frac{\omega_4}{\omega_1}
=(-1)^2\frac{Z_1Z_{2p}Z_{3p}}{Z_{2r}Z_{3r}Z_4}.
\]
Avec $Z_1=30$, $Z_{2p}=17$, $Z_{3p}=80$,
$Z_{2r}=60$, $Z_{3r}=50$ et $Z_4=60$, on obtient
$r\simeq0{,}23$ : la vitesse est divisée par environ quatre.
\end{TLThreeApp}

\subsection{Roue et vis sans fin}

Une vis de $Z_v$ filets entraîne une roue de $Z_r$ dents. Les axes sont
orthogonaux et le rapport vaut en valeur absolue :
\[
\boxed{\left|\frac{\omega_r}{\omega_v}\right|=\frac{Z_v}{Z_r}.}
\]
Un seul étage peut atteindre un rapport de réduction très élevé. En
contrepartie, le rendement est souvent inférieur à $60\,\%$, la vis
subit un effort axial important et l'irréversibilité n'est pas garantie
par la seule géométrie. Pour un faible angle d'hélice, elle peut assurer
une fonction anti-retour utile en levage.

\begin{TLThreeApp}{Réducteur de webcam}
La vis possède un filet et entraîne une roue de $30$ dents. Deux étages
supplémentaires ont les rapports $9/43$ et $9/43$.
\begin{enumerate}
  \item Déterminer le rapport global.
  \item En déduire la vitesse du moteur nécessaire pour obtenir une rotation de la webcam de $90^\circ\,\mathrm{s^{-1}}$.
\end{enumerate}
\end{TLThreeApp}

\subsection{Poulies-courroie et pignons-chaîne}

Ces transmetteurs relient efficacement deux axes parallèles éloignés.
En l'absence de glissement, la vitesse linéaire du brin est la même sur
les deux roues :
\[
R_1\omega_1=R_2\omega_2.
\]
Avec une courroie non croisée ou une chaîne :
\[
\boxed{\frac{\omega_2}{\omega_1}=+\frac{D_1}{D_2}}
\qquad\text{ou}\qquad
\boxed{\frac{\omega_2}{\omega_1}=+\frac{Z_1}{Z_2}.}
\]
Les roues tournent dans le même sens. Une courroie croisée inverse le
sens de rotation.
